\documentclass[11pt]{article}
\usepackage[T1]{fontenc}
\usepackage[margin=1in]{geometry}
\usepackage{amsmath,amssymb,amsthm,mathtools}
\usepackage{graphicx}
\usepackage{microtype}
\usepackage[hidelinks]{hyperref}

\newtheorem{theorem}{Theorem}
\newcommand{\Gn}{G_n}
\newcommand{\gtilde}{\widetilde g}
\newcommand{\Kp}{K_p}

\title{A unit-cube counterexample to the\\discrete covariogram Ball-body inclusion}
\author{Automated Math Discovery Run \texttt{fa\_banach\_001}\\Model: GPT5.6}
\date{22 July 2026}

\begin{document}
\maketitle

\begin{center}
\fbox{\parbox{0.89\linewidth}{\textbf{Claimed status:
counterexample, subject to human review.}  The normalized inclusion marked
unknown in the source fails for $K=[0,1]^n$, $p=1$, and $q=2$ in every
dimension $n\geq1$.}}
\end{center}

\section{Source question}

Alonso-Guti\'errez, Lucas, and Mart\'in Go\~ni ask whether the normalized
inclusion between continuous covariogram Ball bodies, their equation (2.16),
also holds for the Ball bodies of the discrete covariogram
\cite[Section 6, p.~34]{source}.  The source passage is reproduced below.

\begin{figure}[h]
  \centering
  \includegraphics[width=0.95\linewidth]{figures/open_problem_crop.png}
  \caption{The unknown discrete inclusion in the source paper.}
\end{figure}

More explicitly, for every convex body $K\subseteq\mathbb R^n$ with
$0\in K$ and every $0<p<q$, the candidate inclusion is
\begin{equation}\label{eq:candidate}
 \binom{n+q}{n}^{1/q}K_q(\widetilde g_K)
 \stackrel{?}{\subseteq}
 \binom{n+p}{n}^{1/p}K_p(\widetilde g_K).
\end{equation}

\section{Definitions}

Let $G_n(A)=|A\cap\mathbb Z^n|$.  The discrete covariogram of $K$ is
\[
 \widetilde g_K(x)=G_n\bigl(K\cap(x+K)\bigr).
\]
For $p>0$, its $p$-th Ball body is the star set
\[
 K_p(\widetilde g_K)
 =\left\{x\in\mathbb R^n:
 p\int_0^\infty r^{p-1}\widetilde g_K(rx)\,dr
 \geq \widetilde g_K(0)\right\}.
\]
Its radial function satisfies
\begin{equation}\label{eq:radial}
 \rho_{K_p(\widetilde g_K)}(\theta)^p
 =\frac{p}{G_n(K)}\int_0^\infty
 r^{p-1}G_n\bigl(K\cap(r\theta+K)\bigr)\,dr.
\end{equation}

\section{Proof intuition}

The lattice points of the unit cube lie in just two layers perpendicular to
$e_1$.  At zero translation both layers contribute.  As soon as the cube is
translated a positive distance in direction $e_1$, the layer with first
coordinate zero disappears, while the layer with first coordinate one
survives until translation distance one.  Thus the discrete covariogram drops
immediately from $2^n$ to $2^{n-1}$ and then remains constant.  Every radial
$p$-moment therefore has the same $p$-th power, namely $1/2$.  The binomial
normalizations at $p=1$ and $q=2$ then point in the wrong order.

\section{Counterexample}

\begin{theorem}
For every integer $n\geq1$, the convex body $K=[0,1]^n$, with $p=1$ and
$q=2$, violates \eqref{eq:candidate}.  In particular, the proposed discrete
analogue is false even in dimension one.
\end{theorem}

\begin{proof}
The cube contains exactly $2^n$ lattice points, so
$G_n(K)=\widetilde g_K(0)=2^n$.  For $r>0$,
\[
 K\cap(re_1+K)=[r,1]\times[0,1]^{n-1}
 \quad\text{when }0<r\leq1,
\]
and the intersection is empty when $r>1$.  For $0<r\leq1$, its lattice
points are precisely
\[
 \{1\}\times\{0,1\}^{n-1},
\]
so there are $2^{n-1}$ of them.  Hence, for every $p>0$, \eqref{eq:radial}
gives
\begin{equation}\label{eq:half}
 \rho_{K_p(\widetilde g_K)}(e_1)^p
 =\frac{p}{2^n}\int_0^1r^{p-1}2^{n-1}\,dr
 =\frac12.
\end{equation}
The value of the integrand at $r=0$ does not affect the integral.

Now take $p=1$ and $q=2$.  The radial function in direction $e_1$ of the
right-hand side of \eqref{eq:candidate} equals
\[
 \binom{n+1}{n}\rho_{K_1(\widetilde g_K)}(e_1)
 =\frac{n+1}{2}.
\]
The corresponding radial function of the left-hand side is
\[
 \binom{n+2}{n}^{1/2}\rho_{K_2(\widetilde g_K)}(e_1)
 =\frac{\sqrt{(n+1)(n+2)}}{2}.
\]
The latter is strictly larger because $n+2>n+1$.  Inclusion of star sets
would force the opposite inequality between radial functions in every
direction.  Therefore \eqref{eq:candidate} fails.
\end{proof}

\section{Verification, scope, and novelty}

The script \texttt{code/verify\_unit\_cube.py} uses exact rational arithmetic
to enumerate the two lattice layers and check dimensions $1$ through $20$.
It verifies that the squared excess of the proposed left radius over the
right radius is exactly
\[
 \frac{(n+1)(n+2)}4-\frac{(n+1)^2}4=\frac{n+1}{4}>0.
\]
This computation is only a check; the proof above covers all dimensions.

The counterexample addresses exactly the undilated relation marked unknown.
It does not contradict the source's Theorem 6.2, which replaces
$K_p(\widetilde g_K)$ by a dilation of
$K_p(\widetilde g_{K+C_n})$ and includes an additional lattice-count factor.

The run's lightweight indexes had no prior hit.  A bounded search over arXiv,
the published Cambridge version, and exact combinations of the title/id with
``discrete covariogram'', ``Ball bodies'', and ``inclusion'' found no later
resolution or counterexample.  This is not an exhaustive novelty claim.
Human review should check the half-layer count in \eqref{eq:half} and the two
normalization constants.

\begin{thebibliography}{9}
\bibitem{source}
D. Alonso-Guti\'errez, E. Lucas, and J. Mart\'in Go\~ni,
\emph{A discrete approach to Zhang's projection inequality},
Canadian Journal of Mathematics (2025), 1--47;
arXiv:2509.14986.
\url{https://doi.org/10.4153/S0008414X25101624}
\end{thebibliography}

\end{document}

